% Generated by scripts/assemble.py; corrections belong in source/corrections.json.
\examhead{129}{郑州大学 2026 年数学分析试卷}{123}
% p124-b01; source PDF page 124
\blockanchor{p124-b01}
\question{1}(10 分) 求极限 $\lim_{x\to 0^{+}}\ln (2026x)\ln \frac{\ln 2026x}{\ln\frac{x}{2026}}.$\par
% p124-b02; source PDF page 124
\blockanchor{p124-b02}
\question{2}(10 分) 设 $f(x) = \int_{\sqrt{\frac{\pi}{2}}}^{\sqrt{x}} \frac{1}{1 + \cot u^2} \, \mathrm{d}u$, 求积分 $\int_0^{\frac{\pi}{2}} \frac{f(x)}{\sqrt{x}} \, \mathrm{d}x$.\par
% p124-b03; source PDF page 124
\blockanchor{p124-b03}
\question{3}(15 分) 求幂级数 $\sum_{n=2}^{\infty} \frac{x^n}{n^2 - 1}$ 的收敛域及和函数, 并求 $\sum_{n=2}^{\infty} \frac{1}{(n^2 - 1)2^n}$ 的和.\par
% p124-b04; source PDF page 124
\blockanchor{p124-b04}
\question{4}(10 分) 已知曲线积分 $\int_{L} \frac{1}{\varphi(x) + y^2} (x \, \mathrm{d}y - y \, \mathrm{d}x) = A$ (常数), 其中 $\varphi(x)$ 是可导函数且 $\varphi(1) = 1$. $L$ 是绕原点 $(0,0)$ 一周的任意正向闭曲线, 试求 $\varphi(x)$ 及 $A$.\par
% p124-b05; source PDF page 124
\par\Needspace{6\baselineskip}
\blockanchor{p124-b05}
\question{5}(10 分) 设 $z = z(x, y)$ 由方程 $ax + by + cz = F(x^2 + y^2 + z^2)$ 确定, 其中 $F(u)$ 为可导函数, $a, b, c$ 为常数. 证明: 函数 $z = z(x, y)$ 满足方程\par
\nopagebreak[4]
% p124-b06; source PDF page 124
\blockanchor{p124-b06}
\displayformula{(c y - b z) \frac {\partial z}{\partial x} + (a z - c x) \frac {\partial z}{\partial y} = b x - a y.}
% p124-b07; source PDF page 124
\par\Needspace{6\baselineskip}
\blockanchor{p124-b07}
\question{6}(10 分) 设函数 $f(x)$ 在 $[0,1]$ 上具有二阶连续导数, 证明: 在 $(0,1)$ 内存在一点 $\xi$, 使得\par
\nopagebreak[4]
% p124-b08; source PDF page 124
\blockanchor{p124-b08}
\displayformula{\left| f (0) + f (1) - 2 f \left(\frac {1}{2}\right) \right| \leq \frac {1}{4} | f ^ {\prime \prime} (\xi) |.}
% p124-b09; source PDF page 124
\blockanchor{p124-b09}
\question{7}(15 分) 设 $x_0 = \sqrt{6}, x_{n+1} = \sqrt{6 + x_n} (n \geq 0)$.\par
% p124-b10; source PDF page 124
\blockanchor{p124-b10}
(1) 证明: 数列 $\{x_{n}\}$ 收敛, 并求 $\lim_{n\to \infty}x_n$.\par
% p124-b11; source PDF page 124
\blockanchor{p124-b11}
(2) 判断级数 $\sum_{n=1}^{\infty} (-1)^{n-1} \sqrt{3 - x_n}$ 的敛散性, 如果收敛, 判断是绝对收敛还是条件收敛.\par
% p124-b12; source PDF page 124
\blockanchor{p124-b12}
\question{8}(15 分) 设函数列 $\{f_n(x)\}$ 在 $[a, b]$ 上一致收敛于 $f(x)$. 且每个函数 $f_n(x)$ 在 $[a, b]$ 上连续. $\{x_n\} \subset [a, b]$ 满足 $\lim_{n \to \infty} x_n = x_0$, 证明: $\lim_{n \to \infty} f_n(x_n) = f(x_0)$.\par
% p124-b13; source PDF page 124
\blockanchor{p124-b13}
\question{9}(20 分) 设函数 $f(x)$ 在 $[a, b]$ 上可导, 在 $(a, b)$ 内二阶可导. $f(a) = f(b) = 0$, $f_{+}'(a) > 0$, $f_{-}'(b) > 0$.\par
% p124-b14; source PDF page 124
\blockanchor{p124-b14}
(1) 证明: 存在 $c \in (a, b)$, 使得 $f(c) = 0$.\par
% p124-b15; source PDF page 124
\blockanchor{p124-b15}
(2) 证明: 存在 $\xi \in (a, b)$, 使得 $f''(\xi) = 0$.\par
% p124-b16; source PDF page 124
\blockanchor{p124-b16}
(3) 证明: 存在 $\xi_1, \xi_2 \in (a, b)$, 使得 $f''(\xi_1) < 0$, $f''(\xi_2) > 0$.\par
% p124-b17; source PDF page 124
\blockanchor{p124-b17}
\question{10}(15 分) 证明级数 $\sum_{n=1}^{\infty} (-1)^{n-1} \frac{x^n}{n(1 + x^n)}$ 在开区间 $(0,1)$ 上一致收敛于 $f(x)$, 且 $f(x)$ 具有连续导数.\par
% p124-b18; source PDF page 124
\blockanchor{p124-b18}
\question{11}(20 分) 设 $f(x)$ 在 $[0, +\infty)$ 上连续, 常数 $a, b > 0$.\par
% p124-b19; source PDF page 124
\blockanchor{p124-b19}
(1) 若 $\lim_{x\to +\infty}f(x) = k$ ，证明： $\int_0^{+\infty}\frac{f(ax) - f(bx)}{x}\mathrm{d}x = [f(0) - k]\ln \frac{b}{a}.$\par
% p124-b20; source PDF page 124
\blockanchor{p124-b20}
(2) 计算 $\int_0^{+\infty} \frac{\arctan(2x^2) - \arctan(x^2)}{x} \, \mathrm{d}x$.\par
